%% Question bank chapter file - generated by bank_tools.py build (do not edit while the bank is still being built:
%% it is regenerated from the per-test data in Question Bank\data). Every question is one bankq block:
%%   \begin{bankq}[id=..., source={...}, paper=.., qno=.., type=..., subtopic={...}, answer={...}, answer_src=key|solved,
%%                 working={...}, difficulty=1-10, diff_parts={K S I C F}, diff_note={...},
%%                 group=<passage id or empty>, also_in={duplicate copies in other tests}]
%% difficulty = min(10, K+S+I+C+F-1): K concepts 1-3, S steps 1-4, I insight 0-2, C computation 0-2, F format 0-1
%%   ...question LaTeX, options, figures (\bfig / \bfigi -> img/)...
%%   \end{bankq}
%% Passages shared by several questions are bankpassage blocks; their questions carry group=<passage id>.
\documentclass[10pt]{article}
\usepackage[a4paper,margin=1.8cm]{geometry}
\usepackage[T1]{fontenc}
\usepackage{amsmath,amssymb,mathtools}
\usepackage{graphicx}
\usepackage[export]{adjustbox}
\usepackage{xcolor,needspace,keyval}
\usepackage{enumitem}
\usepackage{booktabs,array,tabularx}
\usepackage[hidelinks]{hyperref}
\setlength{\parindent}{0pt}
\setlength{\parskip}{3pt}
\makeatletter
\define@key{bq}{id}{\def\bq@id{#1}}\define@key{bq}{source}{\def\bq@src{#1}}\define@key{bq}{paper}{\def\bq@paper{#1}}
\define@key{bq}{qno}{\def\bq@qno{#1}}\define@key{bq}{type}{\def\bq@type{#1}}\define@key{bq}{subtopic}{\def\bq@sub{#1}}
\define@key{bq}{answer}{\def\bq@ans{#1}}\define@key{bq}{answer_src}{\def\bq@asrc{#1}}\define@key{bq}{working}{\def\bq@work{#1}}
\define@key{bq}{group}{\def\bq@group{#1}}\define@key{bq}{also_in}{\def\bq@also{#1}}
\define@key{bq}{difficulty}{\def\bq@diff{#1}}\define@key{bq}{diff_parts}{\def\bq@dparts{#1}}\define@key{bq}{diff_note}{\def\bq@dnote{#1}}
\newenvironment{bankq}[1][]{\def\bq@work{}\def\bq@also{}\def\bq@group{}\def\bq@diff{?}\def\bq@dparts{}\def\bq@dnote{}\setkeys{bq}{#1}%
  \par\needspace{4\baselineskip}\medskip\noindent{\small\textbf{\bq@id}\hfill\color{gray}[\bq@src, Paper \bq@paper, Q\bq@qno]
  \bq@type\ $\cdot$ \bq@sub\ $\cdot$ \textbf{D\bq@diff}}\par\nopagebreak\relax}
  {\par{\small\color{teal}Answer: \textbf{\bq@ans}\ifx\bq@work\@empty\else\ (solved: \bq@work)\fi
  \ifx\bq@also\@empty\else\par\color{gray}Also in: \bq@also\fi}\par}
\define@key{bp}{id}{\def\bp@id{#1}}\define@key{bp}{source}{\def\bp@src{#1}}\define@key{bp}{questions}{\def\bp@qs{#1}}
\newenvironment{bankpassage}[1][]{\setkeys{bp}{#1}\par\needspace{4\baselineskip}\bigskip\noindent
  {\small\textbf{Passage \bp@id}\hfill\color{gray}[\bp@src; questions \bp@qs]}\par\nopagebreak\relax\itshape}{\par}
\makeatother
\newcommand{\bfig}[2]{\par\smallskip{\centering\includegraphics[width=#1pt,max width=0.62\linewidth,max height=0.32\textheight]{img/#2}\par}\smallskip}
\newcommand{\bfigi}[2]{\includegraphics[width=#1pt,max width=\linewidth,max height=0.32\textheight,valign=c]{img/#2}}
\newcommand{\bmissing}[1]{\textit{[figure missing in the source conversion: \detokenize{#1}]}}
\newenvironment{qoptions}{\begin{itemize}[leftmargin=1.8em,itemsep=1pt,topsep=2pt,label={}]}{\end{itemize}}
\newcommand{\opt}[2]{\item[(#1)] #2}
\providecommand{\fwstart}{}\providecommand{\fwend}{}\providecommand{\qfullsplit}{}
\newcolumntype{L}[1]{>{\raggedright\arraybackslash}p{#1}}
\newcolumntype{C}[1]{>{\centering\arraybackslash}p{#1}}
\newcolumntype{R}[1]{>{\raggedleft\arraybackslash}p{#1}}
\providecommand{\sectionhdr}[1]{}\providecommand{\subsectionhdr}[1]{}


\begin{document}
{\Large\bfseries Differentiation --- Question Bank}\par\smallskip
{\small Mathematics chapter M22. Questions from all converted test papers, grouped by source series. Answers: printed key where it matches the paper, otherwise solved (answer\_src=solved, with working). Exact repeats across papers appear once, with also\_in listing the other copies.}\par\medskip\hrule

\section*{TOAS papers}
\begin{bankq}[id=TOAS06-P1-MATH-S2-Q12, source={TOAS - Test 06 (17 Apr 2022)}, paper=1, qno=12, type=multi, subtopic={Rules}, answer={A,B,C,D}, answer_src=key, difficulty=5, diff_parts={K1 S2 I2 C0 F1}, diff_note={concepts: functional equation with derivative | steps: g = 1 - f is multiplicative; g(x) = x^c; c from f'(1); solve f(x) = 2 | I: transform to a multiplicative function}]
If a derivable function $y=f(x)$ satisfies the condition $f(x)+f(y)=f(x)f(y)+f(xy)\ \forall x,y\in R$ where $f(1)=0$ and $f'(1)=-2$, then the number of solution(s) of the equation $f(x)=2$ is less than or equal to
\begin{qoptions}
\opt{A}{4}
\opt{B}{3}
\opt{C}{2}
\opt{D}{1}
\end{qoptions}
\end{bankq}
\begin{bankq}[id=TOAS06-P1-MATH-S3-Q04, source={TOAS - Test 06 (17 Apr 2022)}, paper=1, qno=4, type=numerical, subtopic={Rules}, answer={1}, answer_src=solved, working={$f(0)=0$ and $|f|\le1$ for $x\to-\infty$ forces the polynomial to be constant, so $f\equiv0$ and $f'(0)=0$: $A=\{0\}$, one integer}, difficulty=4, diff_parts={K1 S1 I2 C0 F1}, diff_note={concepts: polynomial bounded by an exponential expression | steps: f(0) = 0; boundedness as x tends to minus infinity forces f to be constant | I: growth argument}]
$f(x)$ is polynomial which satisfies $|f(x)|\le|e^x-1|\ \forall\ x\in R$ and $A=\{x:x$ is all possible values of $f'(0)\}$, then number of integer in set $A$ is
\end{bankq}
\begin{bankq}[id=TOAS07-P1-MATH-S1-Q01, source={TOAS - Test 07 (01 May 2022)}, paper=1, qno=1, type=single, subtopic={Higher derivatives}, answer={B}, answer_src=key, difficulty=1, diff_parts={K1 S1 I0 C0 F0}, diff_note={concepts: higher derivatives of a Mobius function | steps: y1, y2, y3 in terms of (cx + d); form the ratio}]
If $y=\dfrac{ax+b}{cx+d}$, then value of $\dfrac{y_2^2}{y_1y_3}$, wherever exists, is (where $y_n$ denotes $n^{\mathrm{th}}$ order derivative of differentiable function $y$ and $a,b,c,d$ are constants) -
\begin{qoptions}
\opt{A}{$\dfrac32$}
\opt{B}{$\dfrac23$}
\opt{C}{$\dfrac45$}
\opt{D}{$\dfrac54$}
\end{qoptions}
\end{bankq}
\begin{bankq}[id=TOAS07-P1-MATH-S4-Q09, source={TOAS - Test 07 (01 May 2022)}, paper=1, qno=9, type=numerical, subtopic={Higher derivatives}, answer={4}, answer_src=key, difficulty=7, diff_parts={K2 S3 I1 C1 F1}, diff_note={concepts: derivative of arsinh, differentiating an identity | steps: y'; (1 + x^2) y'^2 = 4y; differentiate -> (1 + x^2) y'' + x y' = 2; differentiate again; read p, q, r, s}, also_in={TOAS - Test 09 (30 May 2022) P2 Q5}]
If $y=\left(\ln\left(x+\sqrt{x^2+1}\right)\right)^2$ and $(x^2+1)\dfrac{d^3y}{dx^3}+(px+q)\dfrac{d^2y}{dx^2}+r\dfrac{dy}{dx}=sy$ , then value of $(p+q+r+s)$ is
\end{bankq}
\begin{bankq}[id=TOAS07-P2-MATH-S1-Q01, source={TOAS - Test 07 (01 May 2022)}, paper=2, qno=1, type=multi, subtopic={Implicit and parametric}, answer={A,B,D}, answer_src=key, difficulty=5, diff_parts={K2 S2 I1 C0 F1}, diff_note={concepts: implicit differentiation, trigonometric equations | steps: find y(0) from the equation at x = 0; differentiate both sides; evaluate y'(0) for each y(0) | I: more than one value of y(0)}]
If $e^{\sin\left(x^2+y^2\right)}=\tan\dfrac{y^2}{4}+\sin^{-1}x$, then $y'(0)$ can be -
\begin{qoptions}
\opt{A}{$\dfrac{1}{3\sqrt\pi}$}
\opt{B}{$-\dfrac{1}{3\sqrt\pi}$}
\opt{C}{$-\dfrac{1}{5\sqrt\pi}$}
\opt{D}{$-\dfrac{1}{3\sqrt{5\pi}}$}
\end{qoptions}
\end{bankq}
\begin{bankq}[id=TOAS07-P2-MATH-S4-Q07, source={TOAS - Test 07 (01 May 2022)}, paper=2, qno=7, type=numerical, subtopic={Higher derivatives}, answer={3}, answer_src=key, difficulty=5, diff_parts={K1 S2 I1 C1 F1}, diff_note={concepts: derivatives of the inverse function | steps: d2x/dy2 = -y''/y'^3; differentiate w.r.t. x and divide by y'; compare coefficients -> k = 3}]
If $\dfrac{d^3x}{dy^3}+\left(\dfrac{dy}{dx}\right)^{-4}\left(\dfrac{d^3y}{dx^3}\right)-k\left(\dfrac{dy}{dx}\right)^{-5}\left(\dfrac{d^2y}{dx^2}\right)^2=0$ is an identity, then the value of $(k)$ is
\end{bankq}
\begin{bankq}[id=TOAS08-P1-MATH-S1-Q04, source={TOAS - Test 08 (15 May 2022)}, paper=1, qno=4, type=multi, subtopic={Implicit and parametric}, answer={A,D}, answer_src=key, difficulty=3, diff_parts={K1 S2 I0 C0 F1}, diff_note={concepts: implicit differentiation of y^x - x^y | steps: find y at x = 1; differentiate both terms; substitute}]
If $y^x-x^y=1$, then $\dfrac{dy}{dx}$ at $x=1$, is -
\begin{qoptions}
\opt{A}{$2(1-\ln2)$}
\opt{B}{$2(1+\ln2)$}
\opt{C}{$\ln(4e^2)$}
\opt{D}{$\ln\left(\dfrac{e^2}{4}\right)$}
\end{qoptions}
\end{bankq}
\begin{bankq}[id=TOAS08-P1-MATH-S1-Q06, source={TOAS - Test 08 (15 May 2022)}, paper=1, qno=6, type=multi, subtopic={Rules}, answer={C}, answer_src=key, difficulty=5, diff_parts={K2 S2 I1 C0 F1}, diff_note={concepts: derivative of a product, binomial sum | steps: f'(0) = (n+1)! times sum of C(n,k)/(k+1); the sum equals (2^(n+1) - 1)/(n+1); solve n | I: sum of C(n,k)/(k+1)}]
If $n\in N$ and $f(x)=({}^nC_0x+1)({}^nC_1x+2)({}^nC_2x+3)({}^nC_3x+4)\ldots({}^nC_nx+n+1)$ and $f'(0)=\dfrac32(7!)$, then $n$ is equal to
\begin{qoptions}
\opt{A}{3}
\opt{B}{4}
\opt{C}{5}
\opt{D}{6}
\end{qoptions}
\end{bankq}
\begin{bankq}[id=TOAS08-P1-MATH-S3-Q06, source={TOAS - Test 08 (15 May 2022)}, paper=1, qno=6, type=numerical, subtopic={Rules}, answer={1.01}, answer_src=solved, working={$P=f'/f$ at 101 $=\sum k(101-k)/(101-k)=\sum k=5050$, $P/5000=1.01$}, difficulty=3, diff_parts={K1 S1 I1 C0 F1}, diff_note={concepts: logarithmic derivative | steps: f'/f = sum of k/(x - k); evaluate at 101 | I: take the logarithm}]
Let $f(x)=(x-1)^{100}(x-2)^{2(99)}(x-3)^{3(98)}\ldots(x-100)^{100}$. If $P=\dfrac{f'(101)}{f(101)}$, then $\dfrac{P}{5000}$ equals
\end{bankq}
\begin{bankq}[id=TOAS08-P2-MATH-S1-Q02, source={TOAS - Test 08 (15 May 2022)}, paper=2, qno=2, type=multi, subtopic={Rules}, answer={D}, answer_src=key, difficulty=3, diff_parts={K1 S1 I0 C1 F1}, diff_note={concepts: derivatives of exponential-log functions | steps: p', r', q' at the points; ratio}]
Let the function $p,q$ and $r$ defined by $p(x)=x^{\ln x}$, $q(x)=(\ln x)^x$, $r(x)=\log_{e^x}x$, then the value of $\dfrac{p'(e^{-1}).r'(e^2)}{q'(e)}$ is (where $y'$ denotes the derivative of $y$)
\begin{qoptions}
\opt{A}{$\dfrac1e$}
\opt{B}{$\dfrac1{e^2}$}
\opt{C}{$\dfrac2e$}
\opt{D}{$\dfrac2{e^2}$}
\end{qoptions}
\end{bankq}
\begin{bankq}[id=TOAS08-P2-MATH-S1-Q04, source={TOAS - Test 08 (15 May 2022)}, paper=2, qno=4, type=multi, subtopic={Higher derivatives}, answer={A,B,C}, answer_src=key, difficulty=4, diff_parts={K1 S2 I1 C0 F1}, diff_note={concepts: second derivative near a singular point | steps: y near pi/2 from sin y = 1 - sin x; y' and y'' behaviour; powers of x | I: local expansion near the vertical tangent}]
For the curve $\sin x+\sin y=1$, $\displaystyle\lim_{x\to0}\left(x^\alpha\right)\dfrac{d^2y}{dx^2}$ exists then $\alpha$ can be
\begin{qoptions}
\opt{A}{$\dfrac32$}
\opt{B}{2}
\opt{C}{$\dfrac52$}
\opt{D}{1}
\end{qoptions}
\end{bankq}
\begin{bankq}[id=TOAS08-P2-MATH-S2-Q09, source={TOAS - Test 08 (15 May 2022)}, paper=2, qno=9, type=single, subtopic={Rules}, answer={A}, answer_src=key, difficulty=1, diff_parts={K1 S1 I0 C0 F0}, diff_note={concepts: derivative matching | steps: compute the four derivatives; match}]
Match List-I with List-II and select the correct answer using the code given below the list.
\begin{center}\footnotesize\renewcommand{\arraystretch}{1.1}\setlength{\extrarowheight}{3.5pt}\setlength{\tabcolsep}{2pt}
\begin{tabular}{|c|l|c|c|}
\hline
 & \multicolumn{1}{c|}{List-I} & & List-II\\[3pt]
\hline
(P) & \parbox[t]{4.4cm}{If $g(x)=f(-x+xf(x))$ and $f(0)=1$, $f'(0)=2$, then value of $g'(0)$ is equal to} & (1) & $\dfrac14$\\[10pt]
(Q) & \parbox[t]{4.4cm}{If $e^y=e^xx^{-y}$, then $\left(\dfrac{dy}{dx}\right)_{x=e}$ is equal to} & (2) & 6\\[8pt]
(R) & \parbox[t]{4.4cm}{Number of points of discontinuity of $y=[2x]$ in $[-1,2]$ is equal to (where $[.]$ denote greatest integer function)} & (3) & 0\\[10pt]
(S) & \parbox[t]{4.4cm}{If $I=\displaystyle\lim_{x\to0}(1+\tan x)^{\frac1{\sin x}}$ then $[I]$ is equal to (where $[.]$ denote greatest integer function)} & (4) & 2\\[10pt]
 & & (5) & 3\\[3pt]
\hline
\end{tabular}
\end{center}
\begin{qoptions}
\opt{A}{$P\to3;Q\to1;R\to2;S\to4$}
\opt{B}{$P\to1;Q\to2;R\to3;S\to5$}
\opt{C}{$P\to2;Q\to1;R\to4;S\to5$}
\opt{D}{$P\to2;Q\to1;R\to3;S\to5$}
\end{qoptions}
\end{bankq}
\begin{bankq}[id=TOAS08-P2-MATH-S2-Q10, source={TOAS - Test 08 (15 May 2022)}, paper=2, qno=10, type=single, subtopic={Implicit and parametric}, answer={B}, answer_src=key, difficulty=2, diff_parts={K1 S1 I1 C0 F0}, diff_note={concepts: implicit curve properties | steps: simplify the inverse trigonometric equation to an algebraic curve; properties for matching | I: combine the three inverse functions}]
If a curve $f(x,y,a,b)=0$ is given by $\sin^{-1}x+\cos^{-1}ay+\cos^{-1}(bxy)=\dfrac\pi2$\\
Match List-I with List-II and select the correct answer using the code given below the list.
\begin{center}\footnotesize\renewcommand{\arraystretch}{1.1}\setlength{\extrarowheight}{3.5pt}\setlength{\tabcolsep}{2pt}
\begin{tabular}{|c|l|c|l|}
\hline
\multicolumn{2}{|c|}{List-I} & \multicolumn{2}{c|}{List-II}\\[3pt]
\hline
(P) & if $a=1$, $b=2$, then & (1) & $(1-x^2)(1-y^2)=x^2y^2$\\[4pt]
(Q) & if $a=2$, $b=1$, then & (2) & $a+b=8$\\[4pt]
(R) & if $a=0$, $b=1$, then & (3) & \parbox[t]{2.6cm}{$\dfrac{dy}{dx}=-\left(\dfrac{y^2+1}{xy}\right)$, $x,y\in R_0$}\\[10pt]
(S) & \parbox[t]{3.1cm}{if $a=b$ and area bounded by rectangle formed by $f(x,y,a,b)=0$ is 1 sq. unit, then} & (4) & $x^2+4y^2=1+3x^2y^2$\\[3pt]
\hline
\end{tabular}
\end{center}
\begin{qoptions}
\opt{A}{$P\to1;Q\to4;R\to2;S\to1$}
\opt{B}{$P\to1;Q\to4;R\to3;S\to2$}
\opt{C}{$P\to4;Q\to1;R\to3;S\to2$}
\opt{D}{$P\to4;Q\to1;R\to3;S\to1$}
\end{qoptions}
\end{bankq}
\begin{bankq}[id=TOAS08-P2-MATH-S3-Q06, source={TOAS - Test 08 (15 May 2022)}, paper=2, qno=6, type=numerical, subtopic={Higher derivatives}, answer={0.20}, answer_src=key, difficulty=3, diff_parts={K1 S1 I1 C0 F1}, diff_note={concepts: Taylor polynomial | steps: express P as a Taylor polynomial about 1; evaluate P(2) and P(3) | I: use Taylor's form about x = 1}]
Let $P(x)$ be a polynomial of degree 4 with $P(1)=1$, $P'(1)=0$, $P''(1)=2$, $P'''(1)=-12$ and $P''''(1)=24$ then value of $\dfrac{P(2)}{P(3)}$ is equal to
\end{bankq}
\begin{bankpassage}[id=TOAS09-P1-MATH-S3-G5-6, source={TOAS - Test 09 (30 May 2022)}, questions={5-6}]
Let $f:[0,\infty)\to[2,\infty)$ be a differentiable increasing and onto function satisfying
\[f^2(x)-2f(x)-\sqrt x=f^2(y)-2f(y)-\sqrt y,\ \forall x,y\in[0,\infty)\]
\end{bankpassage}
\begin{bankq}[id=TOAS09-P1-MATH-S3-Q05, source={TOAS - Test 09 (30 May 2022)}, paper=1, qno=5, type=numerical, subtopic={Rules}, answer={0.50}, answer_src=key, difficulty=2, diff_parts={K1 S1 I0 C0 F1}, diff_note={concepts: limit from the passage | steps: expand f near 0; evaluate the limit}, group=TOAS09-P1-MATH-S3-G5-6]
The value of $\displaystyle\lim_{x\to0^+}\dfrac{f(x)-2}{\sqrt x}$ is equal to -
\end{bankq}
\begin{bankq}[id=TOAS09-P1-MATH-S3-Q06, source={TOAS - Test 09 (30 May 2022)}, paper=1, qno=6, type=numerical, subtopic={Rules}, answer={96.00}, answer_src=key, difficulty=2, diff_parts={K1 S1 I0 C0 F1}, diff_note={concepts: derivative of an inverse function | steps: find x with f(x) = 4; g'(4) = 1/f'(x)}, group=TOAS09-P1-MATH-S3-G5-6]
If $g(x)$ be the inverse of $f(x)$, then $g'(4)$ is equal to -
\end{bankq}
\begin{bankq}[id=TOAS09-P2-MATH-S1-Q19, source={TOAS - Test 09 (30 May 2022)}, paper=2, qno=19, type=single, subtopic={Higher derivatives}, answer={B}, answer_src=key, difficulty=4, diff_parts={K1 S2 I2 C0 F0}, diff_note={concepts: sandwich inequality determines a function | steps: set the left and right inequalities; they force f(x + 4/9) to equal a quadratic; f'' = 9/2 | I: the bounds coincide}]
A real valued function $y=f(x)$ satisfies the relation $f\left(x-\dfrac49\right)+2x\leqslant\dfrac94x^2+\dfrac89\leqslant f\left(x+\dfrac49\right)-2x$. The value of $f''(2)$ is
\begin{qoptions}
\opt{A}{4}
\opt{B}{$\dfrac92$}
\opt{C}{$\dfrac{15}2$}
\opt{D}{$\dfrac{27}2$}
\end{qoptions}
\end{bankq}
\begin{bankq}[id=TOAS09-P2-MATH-S2-Q07, source={TOAS - Test 09 (30 May 2022)}, paper=2, qno=7, type=numerical, subtopic={Rules}, answer={189}, answer_src=solved, working={Derivative of $f^3+g^3+h^3-3fgh$ is zero, constant $=64+125=189$}, difficulty=4, diff_parts={K1 S1 I2 C0 F1}, diff_note={concepts: cyclic derivative system | steps: derivative of f^3 + g^3 + h^3 - 3fgh is zero; evaluate at x = 0 | I: cyclic system gives a constant}]
If $f(x),g(x),h(x)$ be such that $f'(x)=g(x)$, $g'(x)=h(x)$, $h'(x)=f(x)$, $f(0)=4$, $g(0)=5$, $h(0)=0$ then the value of $\left(f^3(x)+g^3(x)+h^3(x)-3f(x)g(x)h(x)\right)$ at $x=2012$ is
\end{bankq}
\end{document}
